\section{An integral class certificate for the conductor-sixteen quartic field}
\label{sec:cubic-conductor16-class-certificate}

Put $p=269$, $t=\sqrt[3]{p}\in\mathbb R$, and
$K=\mathbb Q(\sqrt{-3})$.  Choose
\[
 w=i\sqrt{2+\sqrt2},\qquad F=\mathbb Q(w),\qquad
 L=\mathbb Q(t,w),\qquad M=KL.
 \tag{C16.1}
\]
The choice of the conjugate square root has no effect on these fields.
Let $\psi$ be the character modulo $16$ with
\[
\begin{array}{c|rrrrrrrr}
a&1&3&5&7&9&11&13&15\\\hline
\psi(a)&1&i&i&1&-1&-i&-i&-1,
\end{array}
\qquad \delta=\psi^2=\chi_8.
\tag{C16.2}
\]
All characters are extended by zero on even integers.  This choice
has $\psi(13)=-i$; replacing it by its conjugate merely interchanges
the two quartic factors below.

\begin{theorem}[AC-G142-class: the conductor-sixteen class sector]
\label{AC-G142-class}
For the fields in (C16.1),
\[
                 \operatorname{Cl}(L)[269]=0.
 \tag{C16.3}
\]
The nontrivial cubic-character sector of
$\operatorname{Cl}(M)[269]$ is consequently zero.
The proof gives $h_L=63m$, where $m$ is the index of the five
explicit units below and $269\nmid m$.  It does not assert
$m=1$, a complete class group, or vanishing of the trivial cubic
sector of $\operatorname{Cl}(M)[269]$.
\end{theorem}

\begin{proof}
\proofstep{1. Exact fields and all archimedean factors.}
Writing $z=\zeta_{16}$, one may take $w=z-z^{-1}$ after a
Galois conjugation.  Its stabilizer in
$(\mathbb Z/16\mathbb Z)^\times$ is $\{1,7\}$.
Thus $F/\mathbb Q$ is cyclic quartic with character group
$\{1,\delta,\psi,\bar\psi\}$, and
\[
       w^4+4w^2+2=0,
       \qquad D(F)=8\cdot16^2=2^{11}.
 \tag{C16.4}
\]
Here $\psi$ is primitive since $\psi(9)=-1$, whereas $\delta$
is the primitive quadratic character of conductor eight.
The conductor--discriminant formula proves the discriminant in
(C16.4); in particular the field is unramified at $3$ and $p$.
Its unique quadratic subfield is $\mathbb Q(\sqrt2)$.

The cubic field $T=\mathbb Q(t)$ has discriminant $-3p^2$.
At $p$ the defining polynomial is Eisenstein and tame.  At $3$,
the equation $v+3v^2+3v^3=-30$ has a $\mathbb Z_3$ solution,
by its unit derivative; then $-(1+3v)$ is a cube root of $269$.
The local cubic algebra is therefore
$\mathbb Q_3\times\mathbb Q_3(\sqrt{-3})$.
Its discriminant exponent at three is one, and the polynomial
discriminant excludes all other primes besides three and $p$.
The degrees of $T$ and $F$ are coprime, and their discriminants
are coprime.  At every finite place, one integer algebra is
unramified; its tensor with the other is a product of finite
etale algebras over DVRs and is normal.  The tensor order is
therefore maximal.  It follows that
\[
 [L:\mathbb Q]=12,\quad (r_1,r_2)=(0,6),\quad
 D(L)=2^{33}3^4p^8=19076254851492144304973552812032.
 \tag{C16.5}
\]

The splitting field of $X^3-p$ has group $S_3$ and unique
quadratic subfield $K$.  Its intersection with $F$ is trivial,
since $K$ is ramified at three.  Hence
$\operatorname{Gal}(M/\mathbb Q)=S_3\times C_4$, and $L$ is
fixed by a reflection in the $S_3$ factor.  The normal closure
of this reflection is that entire factor, so every abelian
subextension of $L$ lies in $F$.  All roots of unity in $L$
therefore lie in $F$.  A cyclotomic field contained in $F$ has
degree dividing four.  The nontrivial quadratic possibilities
are excluded by its unique quadratic subfield, and the cyclic
quartic possibilities $\mathbb Q(\zeta_5)=\mathbb Q(\zeta_{10})$
are excluded by ramification.  The remaining degree-four
cyclotomic fields have noncyclic Galois group.  Thus
\[
                         |\mu_L|=2.
 \tag{C16.6}
\]

\proofstep{2. Exact units and a prime-to-269 index.}
Use the ordered rational basis $t^a w^b$, $0\le a<3$, $0\le b<4$,
with index $4a+b$.  The file
\path{exact_units_local_certificate.json} specifies five
rational coefficient arrays and their monic integral degree-twelve
characteristic polynomials, all with constant coefficient $\pm1$.
These are checked in the regular representation: if $T,W$ denote
the multiplication matrices, then
\[
 T^3=269I,\quad W^4+4W^2+2I=0,\quad TW=WT.
 \tag{C16.7}
\]
The raw generator is $x=w-t$.  For each supplied rational
polynomial $q_i(x)$, the verifier computes $q_i(W-T)$ and its
characteristic polynomial exactly over $\mathbb Q$.
Its first column is the asserted tensor array.  Cayley--Hamilton
and the integral constant term show that each value $u_i$ and
its inverse are integral.  Thus these are actual global units,
without any fundamental-unit assumption.
The derived file \path{all_units_exact_certificate.json}
only clears the rational denominators of those identical arrays;
it does not change their values.
For orientation, the first three units are
\[
 \begin{aligned}
 u_1&=3+w^2,\\
 u_2&=\frac{-1423-427w^2+(-218-62w^2)t+(-34-10w^2)t^2}{3},\\
 u_3&=-270+739t-108t^2.
 \end{aligned}
 \tag{C16.8}
\]
The complete two larger arrays, their denominators three, and all
five exact monic polynomials are retained in the finite certificate.

The residue verifier uses the single rational prime
\[
 \ell=206593,\qquad \ell-1=768\cdot269,
 \qquad g=10,\qquad h=g^{768}=49873\pmod\ell.
 \tag{C16.9}
\]
Trial division through $\lfloor\sqrt\ell\rfloor$ proves primality;
the exact power tests for the prime divisors $2,3,269$ of
$\ell-1$ prove that $g$ is primitive.  The three roots of $t^3=p$
are $50819,64197,91577$, and the four roots of $w^4+4w^2+2$
are $39773,79296,127297,166820$, all modulo $\ell$.
Their derivatives are nonzero.  Each of the twelve pairs lifts
by Hensel to an embedding $L\hookrightarrow\mathbb Q_\ell$,
hence gives an actual residue map on the full integer ring.
All unit denominators are invertible modulo $\ell$.
For each such map, define $c_i\in\mathbb F_{269}$ by
$\bar u_i^{768}=h^{c_i}$.
The first five rows, in lexicographic order of the displayed roots, are
\[
 C=\begin{pmatrix}
144&222&85&261&47\\
125&47&85&125&8\\
125&47&85&222&70\\
144&222&85&199&144\\
144&6&214&94&207
\end{pmatrix},\qquad
 \det C=-31248427014\equiv219\pmod{269}.
 \tag{C16.10}
\]
The complete twelve-row matrix is verified by exact modular powers
in \path{269_saturation_certificate.json}.
These residue logarithms kill $269$th powers and $\mu_L$.
Consequently the five units are independent modulo $269$ in the
rank-five lattice $\mathcal O_L^\times/\mu_L$.
They have full rank and their index satisfies
\[
 m=[\mathcal O_L^\times:\langle-1,u_1,\ldots,u_5\rangle]<\infty,
                         \qquad269\nmid m.
 \tag{C16.11}
\]
Indeed a rational relation can be divided to an integer relation
with some coefficient nonzero modulo $269$, contradicting
(C16.10); the same matrix makes the lattice map invertible
modulo $269$, proving the index assertion.

\proofstep{3. The actual regulator and a finite denominator bound.}
Put $\rho=p^{1/3}$ and $\omega=e^{2\pi i/3}$.
The six complex-place representatives are
\[
 (t,w)=(\rho\omega^j,i\sqrt{2+\sqrt2}),\quad
       (\rho\omega^j,i\sqrt{2-\sqrt2}),\qquad j=0,1,2.
 \tag{C16.12}
\]
Form the six-by-five matrix $2\log|\sigma_v(u_i)|$, delete
its last row, and let $R_U$ be the absolute determinant.
This is the projected weighted-log regulator, so $R_U=mR_L$;
there is no extra Euclidean hyperplane-covolume factor.
The exact arrays are evaluated with $3000$-bit outward ball
arithmetic.  Every column sum contains zero, and the determinant
ball is positive and is contained in
\[
                       1309180432<R_U<1309180433.
 \tag{C16.13}
\]
The high precision includes the cancellation in the two large
unit arrays; no provisional regulator is used as an input bound.

For completeness the rational-isolation bound needs no lower
bound for $R_L$.  For degree $n=12$, discriminant $D=D(L)$,
and $r_2=6$, define
\[
 U=\left\lceil\frac{n!4^{r_2}(\lfloor\sqrt D\rfloor+1)}
                              {n^n3^{r_2}}\right\rceil,
 \qquad H=4^n\left\lceil\sqrt[3]{U^4}\right\rceil.
 \tag{C16.14}
\]
Minkowski's ideal-class bound and $\pi>3$ put an integral
representative of every ideal class at norm at most $U$.
For $s=4/3$, Euler factors and an integral comparison give
\[
 h_L\le\#\{\mathfrak a:N\mathfrak a\le U\}
 \le U^{4/3}\zeta_L(4/3)
 \le U^{4/3}\zeta(4/3)^{12}\le4^{12}U^{4/3}\le H.
 \tag{C16.15}
\]
Here there are at most twelve prime ideals over a rational prime,
and $\zeta(4/3)\le1+\int_1^\infty x^{-4/3}dx=4$.
Exact integer arithmetic gives
\[
 U=1318381162569,\qquad H=242535071443530345873408.
 \tag{C16.16}
\]
If $r=h_L/m>1$ and $|r-c|<1/H$ for an integer $c$, then
$m<h_L\le H$.  Unless $h_L=cm$, the integer $h_L-cm$ would
force $|r-c|\ge1/m>1/H$.  Thus the strict interval isolates the
rational value without assuming $m=1$.

\proofstep{4. Artin factorization and finite Dirichlet constants.}
Let $\varrho$ be the standard representation of the $S_3$ splitting
field of $X^3-p$.  The permutation representation for $L$ is
$(1\oplus\varrho)\otimes(1\oplus\delta\oplus\psi\oplus\bar\psi)$.
Including all ramified inertia-invariant factors, Artin formalism gives
\[
 Z:=\operatorname*{res}_{s=1}\zeta_L(s)
 =L(1,\delta)|L(1,\psi)|^2
       L(1,\varrho)L(1,\varrho\delta)|L(1,\varrho\psi)|^2.
 \tag{C16.17}
\]
The quadratic character is primitive modulo eight.
For it the finite Gauss sum is $\sqrt8>0$, as follows by
summing the four elementary eighth-root values.  The finite
Fourier logarithm identity consequently gives
\[
 L(1,\delta)=-\frac1{\sqrt8}\sum_{a=1}^7\delta(a)
          \log(2\sin(\pi a/8))
       =\frac{\log(1+\sqrt2)}{\sqrt2}.
 \tag{C16.18}
\]
The last equality uses $\sin(3\pi/8)/\sin(\pi/8)=1+\sqrt2$.
The identity itself follows by multiplying the absolutely
convergent series $-\log|1-re^{i\theta}|=\sum_{n\ge1}
 r^n\cos(n\theta)/n$ by $\delta(a)$, summing, and then taking
$r\uparrow1$ by Abel's theorem and Dirichlet convergence.
It therefore imports no quadratic class-number or unit-basis claim.
For the primitive odd quartic character, (C16.2) gives exactly
\[
 B_{1,\psi}=\frac1{16}\sum_{a=1}^{15}a\psi(a)=-1-i,
       \qquad |L(1,\psi)|^2=\frac{\pi^2}{8}.
 \tag{C16.19}
\]
The second formula is the primitive odd Dirichlet functional
equation together with the Gauss norm $|\mathfrak g(\psi)|^2=16$.
Only the squared absolute value is used, so no Gauss phase is guessed.

\proofstep{5. All Hecke conductors, including the wild two-part.}
Write $L(s,\varrho)=\sum a_n n^{-s}$.  The coefficients are
multiplicative, with $a_1=1$, $a_{3^v}=1$, and $a_{p^v}=0$
for $v\ge1$.  For primes $\ell\ne3,p$,
\[
 a_\ell=\#\{x\bmod\ell:x^3=p\}-1,\qquad
 a_{\ell^{v+1}}=a_\ell a_{\ell^v}-(-3/\ell)a_{\ell^{v-1}}.
 \tag{C16.20}
\]
At $\ell\equiv2\pmod3$, these are one in even degree and zero
in odd degree.  At $\ell\equiv1\pmod3$, the test
$p^{(\ell-1)/3}$ decides between $a_\ell=2$ and $a_\ell=-1$.
The special factors at three and $p$ follow from the local algebras
already checked in Step 1.  The actual twisted sequences are
$a_n\delta(n)$ and $a_n\psi(n)$; at the new ramified prime two,
$\varrho$ is unramified and its twist has zero inertia invariants.
The finite zero extension therefore inserts the correct local factor.

Write $\varrho=\operatorname{Ind}_{K}^{\mathbb Q}\nu$.
The cubic Kummer character $\nu$ has conductor $(p)$ over $K$:
it is tame and nontrivial at $p$, and $p$ is a cube in $\mathbb Q_3$.
The field $K_2/\mathbb Q_2$ is unramified quadratic.  For $a$ in
its integer ring,
\[
 N(1+4a)\equiv1+4\operatorname{Tr}_{\mathbb F_4/\mathbb F_2}(\bar a)
       \pmod8,\qquad
 N(1+8a)\equiv1+8\operatorname{Tr}_{\mathbb F_4/\mathbb F_2}(\bar a)
       \pmod{16}.
 \tag{C16.21}
\]
The residue trace is onto, $\delta(5)=-1$, and $\psi(9)=-1$.
Their norm pullbacks therefore have exact conductor exponents
three and four respectively.  They are trivial on the next
principal-unit subgroup, and the other local character is
unramified, so no cancellation can lower the conductor.
At $p$ and three the twists are unramified.  Thus the three
primitive Hecke moduli and rational conductors are exactly
\[
 f=p,8p,16p,\qquad
 Q=3f^2=217083,13893312,55573248.
 \tag{C16.22}
\]

\proofstep{6. Primitive finite Fourier transform and theta phase.}
The imaginary quadratic Minkowski bound for $K$ is
$2\sqrt3/\pi<2$, so $h(K)=1$.
Each of the three primitive Hecke characters is represented by
a multiplicative character $\nu_f$ on $\mathcal O_K/f\mathcal O_K$,
zero at nonunits and trivial on the six global units.
For the cubic character the latter also follows from
$9\mid p+1$; for the Dirichlet pullbacks it follows from the norm
of a unit being one.  Put
\[
 a=\frac{2\pi}{\sqrt3f},\qquad
 \Theta(t)=\frac16\sum_{\alpha\in\mathcal O_K}
    \nu_f(\alpha)e^{-a|\alpha|^2t}
       =\sum_{n\ge1}b_ne^{-ant}.
 \tag{C16.23}
\]
The six generators of every ideal give the second identity.

The additive pairing
\[
 (\alpha,\beta)\longmapsto
 \exp\left(2\pi i\operatorname{Tr}_{K/\mathbb Q}
                   \frac{\alpha\bar\beta}{f\sqrt{-3}}\right)
 \tag{C16.24}
\]
is perfect, by the inverse different.  At a unit frequency the
finite Fourier sum is $G\bar\nu_f(\bar\beta)$ by substitution.
At a nonunit frequency at a factor $\mathfrak q^a$ of the modulus,
choose $u\in1+\mathfrak q^{a-1}$ on which the character is
nontrivial, using primitivity; when $a=1$, use a residue unit.
Choose it to be one at the other factors.  Multiplication by $u$
fixes the additive character but multiplies the sum by
$\nu_f(u)\ne1$, so that sum is zero.
Parseval on the $f^2$ residue classes gives $|G|^2=f^2$.
This includes the prime powers eight and sixteen.

The lattice $f\mathcal O_K$ has area $f^2\sqrt3/2$ and
Euclidean dual $2/(f\sqrt{-3})\mathcal O_K$.
For $b=2/\sqrt{-3}$,
$\operatorname{Re}(\alpha\overline{b\beta})
 =\operatorname{Tr}_{K/\mathbb Q}(-\alpha\bar\beta/\sqrt{-3})$.
Thus the two in this dual cancels the one-half in the real-part
formula; no invertibility of two modulo $f$ is asserted.
Poisson summation with the Fourier transform
$\pi(at)^{-1}\exp(-\pi^2|\xi|^2/(at))$ of the Gaussian gives
\[
 \Theta(t)=\varepsilon t^{-1}\overline{\Theta(1/t)},
          \qquad |\varepsilon|=1.
 \tag{C16.25}
\]
Indeed the scalar is $1/(ft)$, the dual exponent is
$a|\beta|^2/t$, and conjugating $\beta$ changes
$\bar\nu_f(\bar\beta)$ to $\bar\nu_f(\beta)$ without changing
its norm.  The scalar Gauss sign is absorbed into $\varepsilon$.
Both constant terms vanish.  Mellin integration, first for
$\operatorname{Re}s>1$ and then by the exponential bounds at both
ends, identifies the integral with $\Gamma(s)a^{-s}L(s)$ at $s=1$.
Splitting at one and substituting $t\mapsto1/t$ gives the exact identity
\[
 L(1)=\sum_{n\ge1}\frac{b_ne^{-an}}n
       +\varepsilon a\sum_{n\ge1}\bar b_n E_1(an),
 \qquad E_1(x)=\int_1^\infty e^{-xt}\frac{dt}{t}.
 \tag{C16.26}
\]
The theta ball at one excludes zero in all three computations.
Therefore its actual phase is
$\varepsilon=\Theta(1)/\overline{\Theta(1)}$, evaluated with
inclusion arithmetic.  It is one for the two real sequences.

\proofstep{7. Tails, the exact rational interval, and descent.}
From (C16.20), $|b_n|\le d(n)\le n$.
Writing $z_a=e^{-a}$ and $H_N=e^{-a(N+1)}$, the omitted tails
for the first sum in (C16.26), the second sum before its factor
$a$, and the theta sum have respective bounds
\[
 \frac{H_N}{1-z_a},\qquad
 \frac{H_N}{a(1-z_a)},\qquad
 \frac{H_N((N+1)-Nz_a)}{(1-z_a)^2}.
 \tag{C16.27}
\]
The middle bound uses $E_1(x)\le e^{-x}/x$.
The script adds each full positive tail radius separately to the
real and imaginary component.  It uses $256$-bit ball arithmetic
with cutoffs $12000,80000,160000$ for (C16.22), and retains the
high-precision regulator ball.  The finite Fourier and closed
quadratic expressions in (C16.18) are also independently evaluated
as balls and their difference contains zero.

The analytic class-number formula in the projected weighted-log
normalization is
\[
 h_LR_L=\frac{2\sqrt{D(L)}}{(2\pi)^6}Z.
 \tag{C16.28}
\]
Combining (C16.13), (C16.17)--(C16.28) gives the strict ball
comparisons
\[
 r=\frac{h_LR_L}{R_U}=\frac{h_L}{m}>1,\qquad
 |r-63|<6\cdot10^{-51},\qquad
 1/H>4.12\cdot10^{-24}.
 \tag{C16.29}
\]
The code compares the balls themselves, not the displayed rounded
bounds.  Its complete output records all weighted-log rows,
conductors, precision, theta and phase balls, tails, final intervals,
and SHA256 hashes of the script and exact unit input.
The selector $63$ enters only in the last comparison, not the
regulator, discriminant, series, or error bounds.
Step 3 now proves $h_L=63m$ exactly.
Together with $269\nmid m$ and $269\nmid63$, this proves (C16.3).

Finally let $s$ be the reflection fixing $L$ inside $M$.
Norm and extension on ideal classes have composites $2$ and $1+s$.
Since $\operatorname{Cl}(L)[269]=0$, the operator $1+s$ vanishes
on the full $269$-primary class group of $M$.
Every conjugate reflection therefore acts as $-1$ there.
Products of reflections generate the cubic rotations, so those
rotations act as one.  This is precisely the claimed vanishing
of the nontrivial cubic sector, with no assertion about the other sector.
\end{proof}

\paragraph{Source and computation scope.}
The retained analytic class-number formula and projected weighted-log
convention are Sutherland, MIT 18.785 Lecture 19 (2019), pp.1 and 4,
\url{https://math.mit.edu/classes/18.785/2019fa/LectureNotes19.pdf}.
The primitive Dirichlet special-value and functional-equation inputs
are DLMF 25.15.3 and 25.15.5--6,
\url{https://dlmf.nist.gov/25.15}.
The classical Minkowski ideal bound is retained as in AC-K19;
the elementary ideal-count isolation argument is proved above.
The finite Fourier cancellation, conductor-sixteen norm calculation,
Poisson normalization, phase, tails, and reflection descent are explicit.
The ball runtime is python-flint 0.9.0 and its documented outward
operations and order-one \texttt{arb.expint}; no implementation proof
is claimed.  The finite unit relations and residue tests use exact
rational and modular arithmetic.
A provisional PARI extraction supplied candidate units and the integer
selector only.  No complete unit basis, provisional class group,
GRH assumption, or unfinished \texttt{bnfcertify} computation is used.
The result supplies only this class sector; its actual strict-group
and growing-window applications require their separate comparisons.
